% ID: SEL-D0003-OB-E02
% ESTADO: cotejada. Fuente y fragmento cotejados de forma independiente; salida compilada revisada.
% Fuente: Selectividad/sel_2019_matematicas_aplicadas/Suplente Junio/Exámenes Matemáticas Aplicadas CCSS - A y B.pdf
% SHA-256: a59474e6895bd41034da53b7060969cbe8cbf7f8272f39fbf16821c95f06c84b
% Página: 2; opción B; ejercicio 2.
% Transcripción de primera pasada basada en la imagen de la página original.
% Sin solución ni modificación de datos o condiciones.

Se considera la función \(f(x)=ax^2+\dfrac{b}{x}\), con \(x\neq 0\), siendo \(a\) y \(b\) dos parámetros reales.
\begin{enumerate}[label=\alph*),leftmargin=2.4em]
  \item \textbf{(1 punto)} Determine el valor de los parámetros \(a\) y \(b\) para que \(f(x)\) tenga un extremo relativo en el punto \((1,\,3)\).
  \item \textbf{(0.75 puntos)} Para \(a=1\) y \(b=2\), razone si en el punto \((1,\,3)\) la función presenta un máximo o un mínimo.
  \item \textbf{(0.75 puntos)} Calcule \(\displaystyle\int\left(x^2+\frac{2}{x}\right)\,dx\).
\end{enumerate}
